% Chapter 22 -- Claim: Transformers Price Stocks Better Than Factor Models
% Requires: amsmath, amssymb, booktabs, array, tabularx, longtable, tcolorbox (breakable, skins), hyperref or url.
% Defines the callout box only if the preamble has not.
\providecolor{fenavy}{HTML}{1B2A6B}
\makeatletter
\@ifundefined{callout}{%
  \newtcolorbox{callout}{breakable, enhanced,
    colback=white, colframe=gray!70, boxrule=0.5pt, arc=1mm,
    borderline west={2.5pt}{0pt}{gray!70},
    left=8pt, right=8pt, top=4pt, bottom=4pt}%
}{}
\makeatother

\providecolor{fenavy}{HTML}{1B2A6B}
\providecolor{feshade}{HTML}{F3F5FA}
\makeatletter
% Shaded box for optional technical material that general readers may skip.
\@ifundefined{technote}{%
  \newtcolorbox{technote}{breakable, enhanced,
    colback=feshade, colframe=fenavy!60, boxrule=0.5pt, arc=1mm,
    left=8pt, right=8pt, top=4pt, bottom=4pt}%
}{}
\makeatother

\chapter{Claim: Transformers Price Stocks Better Than Factor Models}
\label{ch:22}

Anyone who follows the stock market notices that companies do not move alone. When one bank reports trouble, other banks' shares often fall too, and a supplier's bad news can drag down its customers. The newest artificial-intelligence models are designed to learn such connections from data. This chapter asks whether letting a model learn how stocks relate to one another really explains their returns better than the factor models of earlier chapters.

Asset pricing asks why different assets earn different average returns. Since the 1980s its general answer has been the \textbf{stochastic discount factor}\index{stochastic discount factor (SDF)}, a single random variable that prices every asset, which can be thought of as a gauge of how bad economic times are, used to weigh every payoff. The factor models of Chapter~\ref{ch:04} are one way to build it, from a handful of portfolios chosen by researchers. With hundreds of candidate characteristics, later work asked the data to build it instead. Kozak, Nagel and Santosh (2020) worried that a discount factor built from so many portfolios would simply fit noise, so they estimated it from many anomaly portfolios with shrinkage, which pulls weak estimates toward zero. Kelly, Pruitt and Su (2019) started from the observation that a stock's risk changes as its characteristics change, and they let each stock's exposures depend on its characteristics. Chen, Pelger and Zhu (2024) asked whether more flexible, nonlinear relationships could do better still, and they used deep neural networks. Each step priced returns more accurately than the last.

The newest step comes from the technology behind large language models. A \textbf{transformer}\index{transformer|textbf} uses a mechanism called \textbf{attention}\index{attention (neural networks)|textbf}, which lets each word in a sentence look at the other words to decide what it means. Kelly, Kuznetsov, Malamud and Xu (2025) placed a transformer inside an asset-pricing model, so that each stock could ``look at'' other stocks before the model decided how to price it. They reported that this artificial intelligence pricing model priced U.S. stocks with errors about 30\% smaller than the best earlier models. Its tangency portfolio, the portfolio with the highest Sharpe ratio, earned an out-of-sample Sharpe ratio of about four. That is far above anything an investor has earned for long, so the result calls for careful testing.

The stakes are large if the result holds. It would mean that the price of a stock depends on information about other stocks in ways that factor models cannot capture. It would also mean that the methods behind chatbots can improve one of the central models of finance. The claim on trial\index{claims-on-trial method} is: \textbf{an asset-pricing model that places a transformer inside the stochastic discount factor prices the cross-section with much smaller errors than earlier models, and its tangency portfolio earns a higher Sharpe ratio out of sample.}

The lab in Section~\ref{sec:22-lab} replicates the linear-transformer pricing model of Kelly et al.\ (2025) on free anomaly portfolios from the Jensen--Kelly--Pedersen data. It compares the model with the shrinkage baseline of Kozak, Nagel and Santosh (2020). The goal is to check whether attention across assets really prices returns better than a well-built factor model. The original study uses stock-level data from paid databases, so the lab replicates the mechanism one level up, with portfolios as the assets.

By the end of the chapter you will be able to:

\begin{itemize}
  \item explain the stochastic discount factor\index{stochastic discount factor}, and why pricing all assets correctly means finding the portfolio with the highest Sharpe ratio;
  \item describe how shrinkage, IPCA\index{IPCA}, deep networks and adversarial training have been used to estimate it;
  \item compute attention weights by hand, and explain how a transformer shares information across stocks;
  \item judge the claim by out-of-sample pricing errors, Sharpe ratios, results by size group, and costs.
\end{itemize}

\begin{callout}
\textbf{How to read this chapter.} The chapter is written for two kinds of readers at once. If you are new to finance, you can follow the whole argument by reading the main text, the examples and the short \emph{In plain terms} summaries; the box ``For readers new to the stock market'' in Chapter~\ref{ch:01} explains the basic market terms used here. The shaded \emph{Technical note} boxes contain derivations and model formulas, and you can skip them without losing the thread. The Lab and the exercises marked for technical readers are written for readers who want to run the Python code themselves.
\end{callout}

\section{Foundations: the stochastic discount factor and its machine-learning estimators}
\label{sec:22-foundations}

\subsection{The SDF and the pricing equation for excess returns}

An asset-pricing model has one job: to explain why some assets earn higher average returns than others. The most general way to write such a model uses a \textbf{stochastic discount factor}\index{stochastic discount factor|textbf} (SDF). This is a random number $M_{t+1}$ that is high in bad times and low in good times. The model says that every excess return\index{return!excess} $R^e_{i,t+1}$, once multiplied by the SDF, has an expected value of zero:
\begin{equation}
  E_t\!\left[ M_{t+1} R^e_{i,t+1} \right] = 0 \quad \text{for every asset } i.
  \label{eq:22-pricing}
\end{equation}
The intuition is simple. An asset that pays off in bad times, when $M$ is high, is valuable insurance, so investors accept a low average return on it. An asset that pays off in good times must offer a high average return. The CAPM of Chapter~\ref{ch:04} is a special case, in which $M$ falls in a straight line as the market return rises.

Many modern studies write the SDF as one minus the return on a portfolio of the assets:
\begin{equation}
  M_{t+1} = 1 - w_t' R^e_{t+1},
  \label{eq:22-sdf}
\end{equation}
Here $R^e_{t+1}$ is the vector of excess returns, and $w_t$ is the vector of portfolio weights, chosen with information available at time $t$. Estimating an asset-pricing model then means choosing the weights $w_t$. If the weights may depend on each stock's characteristics, such as its size, value and past return, the model can price the cross-section flexibly.

\subsection{From an SDF to the tangency portfolio}

In plain words, finding the discount factor turns out to be the same as finding the portfolio with the best possible Sharpe ratio, the tangency portfolio of Chapter~\ref{ch:12}. The shaded box shows why.

\begin{technote}
\textbf{Technical note: from the SDF to the tangency portfolio.} When the weights are constant, Equation~\eqref{eq:22-pricing} has a familiar solution. Substituting Equation~\eqref{eq:22-sdf} and solving for $w$ gives
\begin{equation}
  w = \Sigma^{-1} \mu,
  \label{eq:22-tangency}
\end{equation}
Here $\mu$ is the vector of expected excess returns and $\Sigma$ is the covariance matrix. (Strictly, second moments replace covariances, a small difference for monthly returns.) These are the weights of the tangency portfolio\index{tangency portfolio} of Chapter~\ref{ch:12}. The SDF portfolio is the portfolio with the highest possible Sharpe ratio, and its squared Sharpe ratio is $\mu'\Sigma^{-1}\mu$.
\end{technote}

This link is why studies of machine-learning SDFs report Sharpe ratios. A model that prices all assets correctly must contain the maximum-Sharpe portfolio. A model that leaves pricing errors behind has missed some return that a better portfolio would capture. Small pricing errors and a high out-of-sample Sharpe ratio are therefore two views of the same success. Studies measure pricing errors in two ways. One is the alphas\index{alpha} of test assets against the SDF portfolio. The other is the \textbf{Hansen--Jagannathan distance}\index{Hansen--Jagannathan distance|textbf}, which summarizes how far the model's SDF is from pricing a set of test assets exactly.

\subsection{Baseline 1: the shrinkage SDF of Kozak, Nagel and Santosh}

With hundreds of candidate factors, Equation~\eqref{eq:22-tangency} cannot be used directly. The sample $\Sigma$ is noisy, its inverse is noisier still, and the weights become extreme. Chapter~\ref{ch:12} called this problem error maximization. Kozak, Nagel and Santosh (2020) shrink the weights with a ridge penalty, exactly as in Chapter~\ref{ch:20}:
\begin{equation}
  \hat w = \left(\hat\Sigma + \gamma I\right)^{-1} \hat\mu,
  \label{eq:22-kns}
\end{equation}
Here $\gamma$ controls how strongly the weights are pulled toward zero, and it is chosen by cross-validation. The penalty is easiest to understand through the principal components\index{principal components} of the factors. It shrinks the weight on each component by the factor $\lambda_j/(\lambda_j+\gamma)$, where $\lambda_j$ is the component's variance. Components with little variance are shrunk the most, because a large average return on them is most likely to be noise. There is also an economic reason. A factor with a high average return but almost no variance would be close to an arbitrage opportunity, and such opportunities should be rare.

Kozak, Nagel and Santosh found that an SDF built from a few characteristic-based factors, such as the four- and five-factor models, could not summarize the cross-section of expected returns well, while an SDF built from a small number of principal components of many factors performed well.

\subsection{Baseline 2: IPCA}

\textbf{Instrumented principal components analysis}\index{IPCA|textbf} (IPCA), introduced by Kelly, Pruitt and Su (2019), lets factor loadings change with a stock's characteristics. The model is
\begin{equation}
  r_{i,t+1} = \beta_{i,t}' f_{t+1} + \varepsilon_{i,t+1}, \qquad \beta_{i,t} = \Gamma' z_{i,t},
  \label{eq:22-ipca}
\end{equation}
Here $f_{t+1}$ is a short vector of unobserved factors, $z_{i,t}$ is the stock's characteristics, and $\Gamma$ is a matrix estimated from the data. A stock's characteristics therefore tell the model how much risk it carries, which is the meaning of the paper's title, ``Characteristics Are Covariances.'' Kelly, Pruitt and Su found that five IPCA factors explained the cross-section of average returns more accurately than existing factor models. The remaining anomaly alphas were small and insignificant, and only about ten characteristics mattered.

\subsection{Deep SDFs: feed-forward networks and adversarial estimation}

A natural next step lets a neural network choose the weights in Equation~\eqref{eq:22-sdf}: $w_{i,t} = g(z_{i,t}, s_t)$. Here $g$ is a network of the kind described in Section~\ref{sec:20-nn}, and $s_t$ is a summary of economic conditions. The network is trained to make the pricing errors as small as possible. This is the same as making the SDF portfolio's Sharpe ratio as large as possible.

Chen, Pelger and Zhu (2024) added a twist from machine learning called \textbf{adversarial training}\index{adversarial training|textbf}. A second network searches for the portfolios that the SDF prices worst. The first network then improves its pricing of exactly those portfolios. The contest continues until the adversary can find no badly mispriced portfolio. They used U.S. stocks from 1967 to 2016, 46 firm characteristics and 178 macroeconomic series, which they summarized with a recurrent network. Their SDF portfolio reached an annual out-of-sample Sharpe ratio of about 2.6, roughly twice that of a network trained only to forecast returns.

\subsection{Attention and the transformer for a panel of stocks}

The models above treat each stock separately: stock $i$'s weight depends only on stock $i$'s own characteristics. Attention removes that restriction. In plain words, each stock looks at every other stock, judges how similar they are, and borrows information from the most similar ones. Give each stock a vector of characteristics $x_i$. Attention computes three transformed versions of it: a \textbf{query} $q_i = W_Q x_i$, a \textbf{key} $k_i = W_K x_i$ and a \textbf{value} $v_i = W_V x_i$. Stock $i$ compares its query with every stock's key, which gives a score for each stock. The scores are turned into weights that sum to one, and stock $i$ takes the weighted average of the other stocks' values:
\begin{equation}
  a_{ij} = \frac{\exp\left(q_i' k_j/\sqrt{d}\right)}{\sum_{l} \exp\left(q_i' k_l/\sqrt{d}\right)}, \qquad h_i = \sum_{j} a_{ij}\, v_j,
  \label{eq:22-attention}
\end{equation}
where $d$ is the length of the query vector. The output $h_i$ is a new description of stock $i$ that includes information from the stocks most similar to it, in the sense that the model has learned.

\textbf{Worked example.} Three stocks have queries and keys of length one, $q = (1, 0, 2)$ and $k = (1, 0, 2)$, and values $v = (0.03, 0.00, -0.02)$. For stock 1, the scores $q_1 k_j$ are $1, 0, 2$. The attention weights are $e^1, e^0, e^2$ divided by their sum, about $0.24, 0.09, 0.67$. Stock 1's new value is $0.24 \times 0.03 + 0.09 \times 0 + 0.67 \times (-0.02) \approx -0.006$. Stock 1 pays most attention to stock 3, whose key is most aligned with its query.

A transformer stacks many such attention steps; the shaded box shows the simplest version used in asset pricing.

\begin{technote}
\textbf{Technical note: the linear transformer.} A \textbf{transformer}\index{transformer|textbf} stacks several such attention layers, each followed by a small feed-forward network, with shortcuts that add each layer's input to its output. A \textbf{linear transformer} drops the exponential and the normalization. In the asset-pricing model of Kelly and coauthors, the linear version gives portfolio weights
\begin{equation}
  w_t = \left(X_t W X_t'\right) X_t \lambda,
  \label{eq:22-linear}
\end{equation}
where $X_t$ stacks all stocks' characteristics in rows, $W$ is a matrix of attention parameters and $\lambda$ a vector of pricing coefficients. The matrix $X_t W X_t'$ has one row and one column for every stock, and its entry $(i,j)$ says how much stock $i$ borrows from stock $j$.
\end{technote}

\noindent\textbf{In plain terms.} An asset-pricing model explains average returns by how much each asset pays off in bad times, and finding the best such model is the same as finding the portfolio with the highest possible Sharpe ratio. Researchers have moved from a few hand-picked factors to statistical shrinkage, to neural networks, and now to transformers that let each stock learn from others.

\section{How it works}
\label{sec:22-how}

\textbf{Why sharing information across stocks might help.} A stock's own characteristics are not the only information about its risk and return. News about a supplier affects its customers. A rise in one bank's funding costs affects other banks. Two stocks with similar characteristics may respond to the same shocks. Earlier models could use such links only if a researcher wrote them in by hand, for example through industry factors. Attention lets the model learn from the data, month by month, which other stocks are informative. Kelly and coauthors checked familiar links: shared analyst coverage, customer--supplier relationships and similar patents. These explained only about 18\% of the variation in their transformer's attention, so most of what the model learned lay outside those links.

\textbf{What the linear surrogate reveals.} Equation~\eqref{eq:22-linear} looks complicated, but it can be rewritten as an ordinary linear model in a very large set of features: every product of three characteristics. With 132 characteristics there are about 2.3 million such products. The linear transformer is therefore a ridge regression of the kind used throughout Chapter~\ref{ch:20}, on an enormous number of interactions. Its weights can be read as signals for timing portfolios built on characteristics. This rewriting is what makes the model, and this chapter's lab, understandable.

\textbf{More parameters than months.} The full transformer has around a million parameters, estimated in rolling windows of 60 months. Classical statistics would expect such a model to overfit. Chapter~\ref{ch:20} described the ``virtue of complexity'' argument, in which very large models with shrinkage can still forecast well. It also described Nagel's (2025) warning that such models may simply reproduce a simple strategy. Kelly, Kuznetsov, Malamud, Xu and Zhang (2024), in ``Large and Deep Factor Models,'' studied deep networks trained to build an SDF. They showed that such a network splits into two parts. One part discovers useful nonlinear characteristics. The other part prices them and behaves like a linear factor model with ridge shrinkage. Their analysis gives a theoretical reason why very large networks need not overfit, and a way to see what they have learned.

\noindent\textbf{In plain terms.} Stocks are linked through suppliers, customers, competitors and shared shocks, and attention lets a model discover such links on its own. Most of what the transformer learned does not match familiar business links, and its huge number of parameters makes careful testing all the more important.

\section{The claim}
\label{sec:22-claim}

Stated so that it can fail:

\begin{callout}
\textbf{An asset-pricing model that uses a transformer to share information across stocks prices the cross-section of U.S. stock returns with smaller out-of-sample pricing errors than factor models and earlier machine-learning models, and the portfolio it implies earns a higher out-of-sample Sharpe ratio, among large stocks and after trading costs.}
\end{callout}

The claim has a \textbf{statistical form}, smaller pricing errors and a higher Sharpe ratio before costs, and an \textbf{economic form}, which requires the advantage to survive among stocks an investor can trade in size and after costs. The tests follow:

\begin{itemize}
  \item \textbf{Out of sample.} Are the gains present in periods not used for training, and after the early decades of the sample?
  \item \textbf{Benchmarks.} Does the transformer beat both simple factor models and strong machine-learning models on the same data?
  \item \textbf{Size.} Are the gains present among mega and large stocks, not only among small and micro stocks?
  \item \textbf{Costs.} How much of the advantage remains after trading costs?
\end{itemize}

\section{The original evidence}
\label{sec:22-evidence}

\subsection{Chen, Pelger and Zhu (2024)}

Chen, Pelger and Zhu\index{Chen, Pelger and Zhu} were among the first to estimate an SDF for individual stocks with deep learning. They trained the network on the pricing condition of Equation~\eqref{eq:22-pricing} rather than on return forecasts, and used an adversary to find the hardest test portfolios. Their SDF beat the benchmarks they considered in Sharpe ratio, explained variation and pricing errors. The most important characteristics were familiar ones, led by past price trends, profitability and capital structure. Conditioning on the economic environment also mattered.

\subsection{Kelly, Kuznetsov, Malamud, Xu and Zhang (2024)}

``Large and Deep Factor Models'' is a theoretical and empirical study of why deep SDF networks work. Its main result is that a deep network trained to build an SDF has two stages. The first discovers nonlinear characteristics. The second prices them like a large linear factor model. An earlier version of the paper reported that, with enough data, out-of-sample performance improved as the network was made deeper. The paper is a working paper and has been substantially revised, so its empirical results should be read in their latest version.

\subsection{Kelly, Kuznetsov, Malamud and Xu (2025)}

Kelly, Kuznetsov, Malamud and Xu\index{AI pricing model (AIPM)|textbf} used monthly U.S. stock returns from 1963 to 2022. The data came from the open-source data of Jensen, Kelly and Pedersen (2023), which Chapter~\ref{ch:19} also uses. They kept 132 of the 153 characteristics in that data, dropped the tiniest stocks, and used the 132 matching factors as test assets. Each model was estimated in rolling 60-month windows, so the out-of-sample returns run from 1968 to 2022. The neural models were averaged over ten random starting points. Table~\ref{tab:22-aipm} shows their main comparison.

\begin{table}[htbp]
\centering
\small
\caption{Out-of-sample performance of SDF models for U.S. stocks, 1968--2022, from Kelly, Kuznetsov, Malamud and Xu (2025), Table 1. The Sharpe ratio is annualized. The pricing error is the Hansen--Jagannathan distance on the 132 test factors; lower is better.}
\label{tab:22-aipm}
\begin{tabular}{@{}lcc@{}}
\toprule
Model & Sharpe ratio & Pricing error \\
\midrule
Fama--French six factors & 1.01 & 0.55 \\
Hou--Xue--Zhang $q$-factors & 1.77 & 0.42 \\
Linear in characteristics (no attention) & 3.60 & 0.15 \\
Shallow random-feature network & 3.87 & 0.13 \\
Linear attention & 3.89 & 0.14 \\
Deep network without attention (MLP) & 4.31 & 0.13 \\
Transformer (AIPM) & 4.57 & 0.09 \\
\bottomrule
\end{tabular}
\end{table}

Four results support the claim. First, the transformer had the highest Sharpe ratio and the smallest pricing error, about 30\% smaller than the deep network without attention. Its alpha against that network was about 1.9\% a year, with a $t$-statistic of 8.4. Second, the advantage held after 2002, when the transformer's Sharpe ratio was 3.37 against 3.07 for the network without attention. Third, adding attention layers helped. The Sharpe ratio rose from about 3.8 with one transformer block to about 4.6 with ten, while making the networks without attention larger did not help. Fourth, the transformer had the highest Sharpe ratio within each size group. Its Sharpe ratio was, however, much lower among mega and large stocks than among small and micro stocks.

\subsection{Questions the papers leave open}

Three questions matter for judging the claim.

\begin{itemize}
  \item \textbf{Costs.} The paper reports no trading costs. Portfolios built from monthly re-estimated weights on thousands of stocks can trade heavily, and Chapter~\ref{ch:17} showed how quickly costs can erase a paper return.
  \item \textbf{Small stocks.} The highest Sharpe ratios come from small and micro stocks, the segment in which Chapter~\ref{ch:19} found anomalies concentrated and Chapter~\ref{ch:20} found machine-learning gains concentrated.
  \item \textbf{Replication.} The paper is a working paper by a single team, and the result has not yet been reproduced independently. Every other model in Table~\ref{tab:22-aipm} also has a Sharpe ratio far above what investors have earned, which is a reminder that these are hypothetical, frictionless returns.
\end{itemize}

\noindent\textbf{In plain terms.} In the main study, the transformer priced stocks more accurately than every earlier model, and its advantage held after 2002 and within each size group. But its returns are hypothetical and before costs, its edge is much smaller among large stocks, and no independent team has yet repeated the result.

\section{Lab~\thechapter: Replicating the linear-transformer SDF on anomaly portfolios}
\label{sec:22-lab}

\begin{callout}
\textbf{Lab~\ref{ch:22} at a glance.} \emph{Papers replicated:} Kelly, Kuznetsov, Malamud and Xu (2025), against the baseline of Kozak, Nagel and Santosh (2020). \emph{Claim on trial\index{claims-on-trial method}:} attention across assets improves an SDF's out-of-sample pricing and Sharpe ratio. \emph{Data:} the free monthly factor returns of Jensen, Kelly and Pedersen (153 U.S. factors, with their theme clusters), and the Fama--French five factors and momentum from the Kenneth R. French Data Library\index{Kenneth French data library}. \emph{Code:} runs as a script or as a notebook, with a prediction prompt before each step. Part D needs PyTorch; the other parts run on a laptop in minutes.
\end{callout}

\textbf{The design.} The original study uses stock-level characteristics, which require paid CRSP and Compustat data. The lab therefore replicates the mechanism one level up. Its ``assets'' are the 153 JKP factors. Each factor's ``characteristics'' are five numbers known at the end of each month: its past 1-, 3- and 12-month returns, its volatility over the past 12 months, and its average return over the past 36 months. Attention then lets each factor borrow information from factors that have behaved similarly. As in the original paper, every model is re-estimated each month on the previous 60 months and judged on the following month.

\begin{itemize}
  \item \textbf{Part A, simple baselines.} The SDF built from the Fama--French five factors plus momentum, and the shrinkage SDF of Equation~\eqref{eq:22-kns} on all 153 factors, with $\gamma$ chosen by cross-validation inside each window.
  \item \textbf{Part B, no attention.} A model in which each factor's weight depends only on its own five characteristics, first linearly and then through a small neural network.
  \item \textbf{Part C, linear attention.} The linear transformer of Equation~\eqref{eq:22-linear}, estimated through its equivalent ridge regression on all products of three characteristics. With five characteristics there are 125 such products, so the regression has a closed-form solution.
  \item \textbf{Part D, a small transformer.} A nonlinear transformer with one to four blocks, trained to maximize the SDF's Sharpe ratio and averaged over ten random starting points.
  \item \textbf{Part E, scoring.} For each model, the out-of-sample Sharpe ratio of the SDF portfolio, its alpha against every other model, and the pricing errors of the 153 factors, reported for the full period and after 2002.
\end{itemize}

\textbf{What readers will find.} Lab Figure~\ref{ch:22}.1 plots the growth of a dollar in each SDF portfolio, scaled to the same volatility. Lab Figure~\ref{ch:22}.2 shows how the transformer's Sharpe ratio changes with the number of blocks. Lab Figure~\ref{ch:22}.3 is a heat map of average attention between factor clusters, showing which groups of factors borrow from which. The test of the claim is whether attention adds a significant alpha over the best model without attention, out of sample and after 2002.

\emph{[Draft note: the figures and numbers are generated by running the lab on data at the time of publication. The book prints the results; the text above does not anticipate them.]}

\subsection{Robustness checks}
\label{sec:22-stress}

The full-sample results can hide weaknesses. The checks below test the claim under the conditions in which it is hardest to meet.

\textbf{1. A placebo.} A flexible model can find structure that is not there. The lab shuffles the characteristics across factors each month, so that a factor's characteristics belong to some other factor. If attention still adds alpha, it is fitting noise. If the advantage disappears, the model is using real information.

\textbf{2. Large stocks.} The lab repeats Part E with JKP factors built from large stocks only, where JKP provide such versions. The original paper's own size-group results show much lower Sharpe ratios among mega and large stocks, and that is the segment an investor can trade in size.

\textbf{3. Costs.} The lab measures how much each model's factor weights change from month to month, and charges an assumed cost for each unit of change. The factors are themselves long--short portfolios, so this understates the true cost. A model that loses its advantage under this charge would therefore lose it in practice too.

\textbf{4. Window length and seeds.} The 60-month window and the network settings are choices. The lab repeats the analysis with 120-month windows and reports the spread of results across random starting points, in the spirit of the robustness grids of Chapter~\ref{ch:24}.

\textbf{5. A simpler explanation.} Following Nagel's (2025) test in Chapter~\ref{ch:20}, the lab compares the transformer's portfolio with a simple factor-timing rule. The rule overweights factors with high recent returns and low volatility. If that rule largely explains the transformer's returns, attention has rediscovered a simple strategy rather than found a new one.

\subsection*{Your verdict}

Reach your own verdict on the claim from the numbers in your own run, using the scale of Section~\ref{sec:01-method}. Then compare it with the published evidence summarized in Section~\ref{sec:22-verdict}. If your verdict differs, explain which of your results led you there.

\section{Where the evidence stands}
\label{sec:22-verdict}

The published studies cited in this chapter point to an early picture: \textbf{one careful study found that a transformer prices U.S. stocks better than factor models and earlier machine-learning models out of sample. Its value after costs and among large stocks has not been tested, and no one has yet reproduced the result independently.} The table summarizes where the evidence supports the claim and where it does not, as of November 2026, and names the source of each finding. The lab lets you verify these findings with free data.

{\small
\begin{longtable}{@{}>{\raggedright\arraybackslash}p{0.19\linewidth}>{\raggedright\arraybackslash}p{0.10\linewidth}>{\raggedright\arraybackslash}p{\dimexpr 0.48\linewidth-6\tabcolsep\relax}>{\raggedright\arraybackslash}p{0.23\linewidth}@{}}
\toprule
Situation & Holds? & Why & Evidence \\
\midrule
\endfirsthead
\multicolumn{4}{@{}l}{\emph{(continued)}} \\
\toprule
Situation & Holds? & Why & Evidence \\
\midrule
\endhead
\midrule
\multicolumn{4}{r@{}}{\emph{(continued on the next page)}} \\
\endfoot
\bottomrule
\endlastfoot
Against factor models & \textbf{Yes} & Far higher Sharpe ratio and far smaller pricing errors than the Fama--French six factors and the $q$-factors. & Kelly, Kuznetsov, Malamud and Xu (2025) \\
\addlinespace[3pt]
Against networks without attention & \textbf{Yes} & Pricing errors about 30\% smaller than the best deep network without attention, with an alpha of about 1.9\% a year ($t = 8.4$). & Kelly, Kuznetsov, Malamud and Xu (2025) \\
\addlinespace[3pt]
Deep SDFs in general & \textbf{Yes} & A network trained on the pricing condition beat its benchmarks, with familiar characteristics doing most of the work. & Chen, Pelger and Zhu (2024) \\
\addlinespace[3pt]
After 2002 & \textbf{Partly} & The advantage was smaller but still present: a Sharpe ratio of 3.37 against 3.07. & Kelly, Kuznetsov, Malamud and Xu (2025) \\
\addlinespace[3pt]
Large stocks & \textbf{Weakly} & The transformer was best within each size group, but its Sharpe ratio was much lower among mega and large stocks. & Kelly, Kuznetsov, Malamud and Xu (2025) \\
\addlinespace[3pt]
After trading costs & \textbf{Untested} & No costs were reported. & None yet \\
\addlinespace[3pt]
Independent replication & \textbf{Untested} & The result comes from a working paper by a single team. & None yet \\
\end{longtable}
}

Every model in Table~\ref{tab:22-aipm}, not only the transformer, reports a Sharpe ratio far above what investors have earned, a reminder that these are hypothetical, frictionless returns.

\textbf{What would change the picture.} The economic form of the claim would gain support if a cost-aware version, judged among large stocks, kept its advantage and an independent team reproduced it. It would lose support if the transformer's returns were largely explained by a simple factor-timing rule.

\subsection{Open questions}
\label{sec:22-open}

\begin{itemize}
  \item \textbf{What does attention learn?} Observable links between firms explain only a small part of the transformer's attention. What does the rest represent?
  \item \textbf{Do the gains survive costs?} A cost-aware version of the model, trained to maximize net rather than gross Sharpe ratio, has not yet been published.
  \item \textbf{Is it risk or mispricing?} An SDF with a Sharpe ratio above four implies extremely volatile marginal utility if it reflects risk. If it reflects mispricing, how long can it last once the model is widely used?
  \item \textbf{Beyond U.S. stocks.} Does attention help in other countries, where the JKP data also allow tests, and in other asset classes?
\end{itemize}

\section{What practitioners and investors can use}
\label{sec:22-investors}

The evidence in Section~\ref{sec:22-verdict} suggests several practical lessons.

\textbf{Read Sharpe ratios above three as research results, not as returns.} Every machine-learning SDF in this chapter reports Sharpe ratios that no fund has delivered for long. They are hypothetical, frictionless and long--short, and they include stocks that are hard to trade.

\textbf{The lasting lesson is about information.} Stocks are not independent. The attention result suggests that information about related firms helps to explain returns. An investor can use that idea without a transformer, for example by watching how a company's suppliers, customers and competitors are moving.

\textbf{Treat ``AI-powered'' as a claim to test.} A product built on these models should be judged by the same questions as any other strategy.

\textbf{Questions to ask a manager who uses AI to price stocks:}

\begin{enumerate}
  \item What is the strategy's Sharpe ratio after costs, among stocks you can trade in your size?
  \item How much of the return comes from small and micro stocks?
  \item How does the model compare with a simple linear model on the same characteristics, and with a simple factor-timing rule?
  \item How often is the model re-estimated, how much does the portfolio turn over, and what happened in periods when it was not re-trained?
  \item Has anyone outside your team reproduced the results?
\end{enumerate}

\begin{callout}
\textbf{What Failed: the quant quake of August 2007.} In the first two weeks of August 2007, many quantitative equity funds suffered large and sudden losses, followed by a partial rebound. These funds held long--short portfolios built from similar characteristics, such as value, momentum and quality. Khandani and Lo (2011) traced the episode to the rapid sale of one or more large portfolios. The sale pushed down the prices of the stocks that all these funds held long, and pushed up the prices of the stocks they held short. The models were not wrong about average returns. They were wrong about risk. Funds that thought they held different strategies had built nearly the same portfolio from the same characteristics. The claim that a characteristic-based model's backtested risk describes its live risk \textbf{fails} when many investors use similar models. A model that learns from the same public characteristics as everyone else may end up in the same crowded trades.
\end{callout}

\section*{Exercises}
\addcontentsline{toc}{section}{Exercises}

\subsection*{Conceptual (all readers)}

\begin{enumerate}
  \item Explain in words why an asset that pays off in bad times should earn a lower average return, using the SDF of Equation~\eqref{eq:22-pricing}.
  \item Show that if $M = 1 - w'R^e$ prices two uncorrelated assets with expected excess returns of 4\% and 6\% and variances of 0.04 and 0.09, the weights are $w = (1, 0.67)$. What is the Sharpe ratio of the SDF portfolio?
  \item In the shrinkage SDF of Equation~\eqref{eq:22-kns}, a principal component has variance 0.0001 and another has variance 0.01, and $\gamma = 0.001$. By what fraction is each component's weight shrunk? Why is this economically sensible?
  \item Repeat the attention worked example for stock 3. Which stock does stock 3 pay most attention to, and what is its new value?
  \item Explain why a Sharpe ratio of 4.57 in a research paper does not mean that an investor could have earned it.
  \item A news headline says that ``AI now understands the stock market better than economists do.'' Without any calculation, use this chapter to explain what the evidence actually shows, and what has not yet been tested.
\end{enumerate}

\subsection*{Lab extensions (technical readers)}

\begin{enumerate}
  \setcounter{enumi}{6}
  \item In Part C, add the JKP theme cluster of each factor as extra characteristics. Does attention concentrate within clusters?
  \item Run the placebo of Robustness check 1 one hundred times. Where does the real model's alpha fall in the distribution of placebo alphas?
  \item In Part D, plot the out-of-sample Sharpe ratio against the number of transformer blocks for windows of 60 and 120 months. Does depth help in both?
  \item Charge costs of 5, 10 and 20 basis points per unit change in factor weights. At what cost does the transformer's advantage over the shrinkage SDF disappear?
  \item Repeat the full lab with the JKP factors for developed markets outside the United States. Is the ranking of models the same?
\end{enumerate}
